%============================================================
% Chapter 16: Walsh--Hadamard and Bitwise Convolutions
% STATUS: PLANNED  (generated by tools/outline.py; replace this file when the chapter is written)
%============================================================
\chapter{Walsh--Hadamard and Bitwise Convolutions}
\label{ch:fwht}

\begin{tcolorbox}[enhanced, colback=white, colframe=secondaryblue, boxrule=1pt, arc=2mm]
\textbf{\large\textcolor{primaryblue}{Convolve two arrays indexed by masks with XOR, AND or OR in $\bigO{n \log n}$: the Walsh--Hadamard, zeta and M\"obius transforms.}}
\end{tcolorbox}

\chapterfacts{2200}{3200}{4}{8}{30}

\begin{plannedbox}[Planned Contents --- this chapter is not written yet]
\textbf{Part A --- Lessons}
\begin{enumerate}[leftmargin=*, label=\textbf{A\arabic*.}, itemsep=1pt]
    \item XOR convolution: the Walsh--Hadamard transform and its inverse
    \item AND and OR convolutions: superset and subset sums as transforms (SOS DP is the DP book)
    \item Working modulo a prime, and the division by $n$
    \item Counting pairs and triples with a given XOR over huge multiplicities
    \item Subset convolution is the DP book; this chapter owns the three bitwise transforms
\end{enumerate}
\textbf{Part B --- Patterns} (with candidate problems, checked when the chapter is written)
\begin{enumerate}[leftmargin=*, label=\textbf{P\arabic*.}, itemsep=1pt]
    \item \textbf{XOR convolution} --- Bitwise Xor Convolution {\small\color{darkgray}(Library Checker)}
    \item \textbf{AND and OR convolutions} --- Bitwise And Convolution {\small\color{darkgray}(Library Checker)}
    \item \textbf{Transforms in counting problems} --- Triple {\small\color{darkgray}(CF 1119H)}
    \item \textbf{Powers of a transform}
\end{enumerate}
\end{plannedbox}
